% aoa-arith.tex -- the arithmetic commands.
%   \add{a}{b}  or  \add{a,b,c,...}   (any number of rows)
%   \sub{a}{b}      \mul{a}{b}      \dvd{n}{d}  (n divided by d)
% STAND-INS: this is the one file to replace (or turn into thin wrappers)
% when Alex's own package is ready. Problem files don't change.
%
% Answers are computed exactly, to any number of digits (bigintcalc for
% + - x, integer long division for \dvd).
\makeatletter
\RequirePackage{bigintcalc,array}

% \checkanswers: write every computed answer to <jobname>.ans, for checking
\newwrite\aoa@ansfile \newif\ifaoa@logging
\newcommand{\checkanswers}{\aoa@loggingtrue\immediate\openout\aoa@ansfile=\jobname.ans}
\def\aoa@logans#1#2{\ifaoa@logging
  \immediate\write\aoa@ansfile{#1|#2|\expandafter\strip@prefix\meaning\aoa@ans}\fi}

% ------------------------------------------------------------ commands
\newcommand{\add}[1]{\@ifnextchar\bgroup{\aoa@addtwo{#1}}{\aoa@arith{+}{add}{#1}}}
\def\aoa@addtwo#1#2{\aoa@arith{+}{add}{#1,#2}}
\newcommand{\sub}[2]{\aoa@arith{-}{sub}{#1,#2}}
\newcommand{\mul}[2]{\aoa@arith{\times}{mul}{#1,#2}}
\newcommand{\dvd}[2]{\aoa@division{#1}{#2}}

% #1 printed operator, #2 add|sub|mul, #3 comma list of operands
% (spaces and line breaks in the list are removed first)
\def\aoa@arith#1#2#3{%
  \edef\aoa@L{\zap@space#3 \@empty}%
  \expandafter\aoa@@arith\expandafter{\aoa@L}{#1}{#2}}
\def\aoa@@arith#1#2#3{\aoa@@@arith{#2}{#3}{#1}}
\def\aoa@@@arith#1#2#3{%
  \aoa@compute{#2}{#3}%
  \global\let\aoa@lastans\aoa@ans   % for the answer key
  \aoa@logans{#2}{#3}%
  \def\aoa@tmp{compact}%
  \ifx\aoa@style\aoa@tmp
    \aoa@compact{#1}{#3}%
  \else
    \aoa@stacked{#1}{#3}%
  \fi}

\def\aoa@show{show}\def\aoa@solution{solution}\def\aoa@space{space}

% compact: 4807 + 2971 (= 7778 with answers=show)
\def\aoa@compact#1#2{%
  \ifx\aoa@answers\aoa@show
    $\aoa@inline{#1}{#2} = \aoa@ans$%
  % solutions and workbook: the problem appears once, in the solution's form
  \else\ifx\aoa@answers\aoa@solution
    \aoa@solutionbox{#1}{#2}%
  \else\ifx\aoa@answers\aoa@space
    \aoa@stack{#1}{#2}{\phantom}%
  \else
    $\aoa@inline{#1}{#2}$%
  \fi\fi\fi}
\def\aoa@inline#1#2{%
  \def\aoa@sep{}%
  \@for\aoa@x:=#2\do{\aoa@sep\aoa@x\def\aoa@sep{#1}}}

\def\aoa@stacked#1#2{%
  \ifx\aoa@answers\aoa@show
    \aoa@stack{#1}{#2}{\@firstofone}%
  \else\ifx\aoa@answers\aoa@solution
    \aoa@solutionbox{#1}{#2}%
  \else\ifx\aoa@answers\aoa@space
    \aoa@stack{#1}{#2}{\phantom}%
  \else
    \aoa@stack{#1}{#2}{}%
  \fi\fi\fi}

% Stand-in for the worked solution: the stacked form, answer in color.
\def\aoa@solcolor{\textcolor{blue!60!black}}
\def\aoa@solutionbox#1#2{\aoa@stack{#1}{#2}{\aoa@solcolor}}

% Stacked layout: integer parts right-aligned, decimal parts left-aligned,
% so decimal points line up; the operator sits on the last row.
% #3 is how to print the answer row: empty (no answer), \@firstofone,
% \phantom (room only) or \aoa@solcolor.
\newtoks\aoa@rows
\newcount\aoa@n \newcount\aoa@i
\def\aoa@stack#1#2#3{%
  \aoa@n=0 \@for\aoa@x:=#2\do{\advance\aoa@n 1 }%
  \aoa@rows={}\aoa@i=0
  \@for\aoa@x:=#2\do{%
    \advance\aoa@i 1
    \expandafter\aoa@split\expandafter{\aoa@x}%
    \edef\aoa@row{\ifnum\aoa@i=\aoa@n\noexpand#1\fi
      & \aoa@int & \aoa@frac \noexpand\\}%
    \aoa@rows=\expandafter{\the\expandafter\aoa@rows\aoa@row}}%
  \expandafter\aoa@split\expandafter{\aoa@ans}%
  \def\aoa@tmp{#3}%
  $\begin{array}[t]{@{}c@{\,}r@{}l@{}}
     \the\aoa@rows \hline
     \ifx\aoa@tmp\@empty
       & \vphantom{0} &
     \else
       & #3{\aoa@int} & #3{\aoa@frac}
     \fi
   \end{array}$}

% \aoa@split{12.34} -> \aoa@int=12, \aoa@frac=.34 (empty if no point)
\def\aoa@split#1{\aoa@@split#1..\@nil}
\def\aoa@@split#1.#2.#3\@nil{%
  \def\aoa@int{#1}%
  \def\aoa@tmpb{#2}%
  \ifx\aoa@tmpb\@empty\def\aoa@frac{}\else\def\aoa@frac{.#2}\fi}

% ------------------------------------------------ exact + - x on decimals
% Each operand is scaled to an integer, computed with bigintcalc, and the
% decimal point put back.
\newcount\aoa@cnt \newcount\aoa@K \newcount\aoa@k
% \aoa@parse{x}: \aoa@I integer digits, \aoa@F fraction digits, \aoa@k their count
\def\aoa@parse#1{\aoa@@parse#1..\@nil}
\def\aoa@@parse#1.#2.#3\@nil{%
  \def\aoa@I{#1}\ifx\aoa@I\@empty\def\aoa@I{0}\fi
  \def\aoa@F{#2}\aoa@len{#2}\aoa@k=\aoa@cnt}
\def\aoa@len#1{\aoa@cnt=0 \aoa@@len#1\aoa@end}
\def\aoa@@len#1{\ifx#1\aoa@end\else\advance\aoa@cnt 1 \expandafter\aoa@@len\fi}
\def\aoa@end{\aoa@end}
% \aoa@zeros{n}: n zeros (expandable)
\def\aoa@zeros#1{\ifnum#1>0 0\expandafter\aoa@zeros\expandafter{\the\numexpr#1-1\relax}\fi}
% \aoa@scaled{x}{K} -> \aoa@S = x * 10^K as a normalized integer
\def\aoa@scaled#1#2{%
  \aoa@parse{#1}%
  \edef\aoa@S{\aoa@I\aoa@F\aoa@zeros{\numexpr#2-\aoa@k\relax}}%
  \edef\aoa@S{\bigintcalcNum{\aoa@S}}}

\def\aoa@compute#1#2{%
  \def\aoa@tmp{#1}\def\aoa@tmpm{mul}%
  \ifx\aoa@tmp\aoa@tmpm
    % multiply: decimals add up
    \aoa@K=0 \def\aoa@acc{1}%
    \@for\aoa@x:=#2\do{%
      \expandafter\aoa@parse\expandafter{\aoa@x}%
      \expandafter\aoa@scaled\expandafter{\aoa@x}{\aoa@k}%
      \advance\aoa@K\aoa@k
      \edef\aoa@acc{\bigintcalcMul{\aoa@acc}{\aoa@S}}}%
    \aoa@unscale{\aoa@acc}{\aoa@K}\aoa@stripzeros
  \else
    % add or subtract: work at the largest number of decimals
    \aoa@K=0
    \@for\aoa@x:=#2\do{%
      \expandafter\aoa@parse\expandafter{\aoa@x}%
      \ifnum\aoa@k>\aoa@K \aoa@K=\aoa@k \fi}%
    \let\aoa@acc\relax
    \@for\aoa@x:=#2\do{%
      \expandafter\aoa@scaled\expandafter{\aoa@x}{\aoa@K}%
      \ifx\aoa@acc\relax \let\aoa@acc\aoa@S
      \else
        \def\aoa@tmpm{sub}%
        \ifx\aoa@tmp\aoa@tmpm
          \edef\aoa@acc{\bigintcalcSub{\aoa@acc}{\aoa@S}}%
        \else
          \edef\aoa@acc{\bigintcalcAdd{\aoa@acc}{\aoa@S}}%
        \fi
      \fi}%
    \aoa@unscale{\aoa@acc}{\aoa@K}%
  \fi}

% \aoa@unscale{R}{K}: \aoa@ans = R / 10^K written with K decimals
\def\aoa@unscale#1#2{%
  \edef\aoa@R{#1}\def\aoa@sign{}%
  \expandafter\aoa@checksign\aoa@R\@nil
  \ifnum#2=0 \edef\aoa@ans{\aoa@sign\aoa@R}%
  \else
    \expandafter\aoa@len\expandafter{\aoa@R}%
    \ifnum\aoa@cnt>#2 \else
      \edef\aoa@R{\aoa@zeros{\numexpr#2-\aoa@cnt+1\relax}\aoa@R}%
      \expandafter\aoa@len\expandafter{\aoa@R}%
    \fi
    \aoa@i=\numexpr\aoa@cnt-#2\relax
    \def\aoa@head{}\def\aoa@tail{}%
    \expandafter\aoa@splitat\aoa@R\aoa@end
    \edef\aoa@ans{\aoa@sign\aoa@head.\aoa@tail}%
  \fi}
\def\aoa@checksign#1#2\@nil{\ifx-#1\def\aoa@sign{-}\def\aoa@R{#2}\fi}
% move the first \aoa@i digits to \aoa@head, the rest to \aoa@tail
\def\aoa@splitat#1{%
  \ifx#1\aoa@end\else
    \ifnum\aoa@i>0 \edef\aoa@head{\aoa@head#1}\advance\aoa@i -1
    \else \edef\aoa@tail{\aoa@tail#1}\fi
    \expandafter\aoa@splitat\fi}
% drop trailing zeros after the point (and the point itself if nothing is left)
\def\aoa@stripzeros{%
  \expandafter\aoa@split\expandafter{\aoa@ans}%
  \ifx\aoa@frac\@empty\else
    \edef\aoa@F{\expandafter\@gobble\aoa@frac}%
    \aoa@revstrip
    \ifx\aoa@F\@empty \edef\aoa@ans{\aoa@int}%
    \else \edef\aoa@ans{\aoa@int.\aoa@F}\fi
  \fi}
\def\aoa@revstrip{%
  \edef\aoa@last{\expandafter\aoa@lastdigit\aoa@F\aoa@end}%
  \if0\aoa@last
    \edef\aoa@F{\expandafter\aoa@dropl\aoa@F\aoa@end}%
    \ifx\aoa@F\@empty\else\expandafter\expandafter\expandafter\aoa@revstrip\fi
  \fi}
\def\aoa@lastdigit#1#2{\ifx#2\aoa@end #1\else\expandafter\aoa@lastdigit\expandafter#2\fi}
\def\aoa@dropl#1#2{\ifx#2\aoa@end\else #1\expandafter\aoa@dropl\expandafter#2\fi}

% ------------------------------------------------------------- division
% n / d for whole numbers: the decimal, with a bar over the repeating
% part. Long repeats (more than 12 digits) show the first 24 digits.
\newcount\aoa@q \newcount\aoa@r \newcount\aoa@pos \newcount\aoa@stamp
\def\aoa@division#1#2{%
  \aoa@longdiv{#1}{#2}%
  \global\let\aoa@lastans\aoa@ans
  \aoa@logans{dvd}{#1,#2}%
  \def\aoa@tmp{compact}%
  \ifx\aoa@answers\aoa@show
    \ifx\aoa@style\aoa@tmp $#1 \div #2 = \aoa@ans$%
    \else \aoa@divstack{#1}{#2}{\@firstofone}\fi
  \else\ifx\aoa@answers\aoa@solution \aoa@divstack{#1}{#2}{\aoa@solcolor}%
  \else\ifx\aoa@answers\aoa@space \aoa@divstack{#1}{#2}{\phantom}%
  \else
    \ifx\aoa@style\aoa@tmp $#1 \div #2$%
    \else \aoa@divstack{#1}{#2}{}\fi
  \fi\fi\fi}
% long-division layout: answer over the bar, divisor ) dividend
\def\aoa@divstack#1#2#3{%
  \def\aoa@tmp{#3}%
  $\begin{array}[t]{@{}r@{}l@{}}
     & \ifx\aoa@tmp\@empty\vphantom{0}\else\;#3{\aoa@ans}\fi \\ \cline{2-2}
     #2\,\big) & \;#1
   \end{array}$}

\def\aoa@longdiv#1#2{%
  \global\advance\aoa@stamp 1
  \aoa@q=\numexpr#1/#2\relax
  \ifnum\numexpr\aoa@q*#2\relax>#1 \advance\aoa@q -1 \fi
  \aoa@r=\numexpr#1-\aoa@q*#2\relax
  \edef\aoa@whole{\the\aoa@q}%
  \def\aoa@digits{}\aoa@pos=0 \let\aoa@reppos\relax
  \loop
    \ifnum\aoa@r=0 \aoa@pos=-1 \fi
    \ifnum\aoa@pos>-1
      \expandafter\ifx\csname aoa@seen@\the\aoa@stamp @\the\aoa@r\endcsname\relax
        \expandafter\edef\csname aoa@seen@\the\aoa@stamp @\the\aoa@r\endcsname{\the\aoa@pos}%
        \aoa@q=\numexpr(10*\aoa@r)/#2\relax
        \ifnum\numexpr\aoa@q*#2\relax>\numexpr10*\aoa@r\relax \advance\aoa@q -1 \fi
        \aoa@r=\numexpr10*\aoa@r-\aoa@q*#2\relax
        \edef\aoa@digits{\aoa@digits\the\aoa@q}%
        \advance\aoa@pos 1
      \else
        \edef\aoa@reppos{\csname aoa@seen@\the\aoa@stamp @\the\aoa@r\endcsname}%
        \aoa@pos=-1
      \fi
    \fi
  \ifnum\aoa@pos>-1 \repeat
  \ifx\aoa@digits\@empty
    \let\aoa@ans\aoa@whole          % divides evenly
  \else\ifx\aoa@reppos\relax
    \edef\aoa@ans{\aoa@whole.\aoa@digits}% terminates
  \else
    % split digits at the start of the repeat
    \aoa@i=\aoa@reppos\relax
    \def\aoa@head{}\def\aoa@tail{}%
    \expandafter\aoa@splitat\aoa@digits\aoa@end
    \expandafter\aoa@len\expandafter{\aoa@tail}%
    \ifnum\aoa@cnt>12
      \aoa@i=24 \edef\aoa@R{\aoa@head\aoa@tail\aoa@tail}%
      \let\aoa@hh\aoa@head
      \def\aoa@head{}\def\aoa@tail{}%
      \expandafter\aoa@splitat\aoa@R\aoa@end
      \edef\aoa@ans{\aoa@whole.\aoa@head\noexpand\ldots}%
    \else
      \edef\aoa@ans{\aoa@whole.\aoa@head\noexpand\overline{\aoa@tail}}%
    \fi
  \fi\fi}
\makeatother
